% ======================================================================
% sections/limitations_future.tex
%
% Limitations and Future Work scaffold for the paper:
%   Triple Humanoid 24/7 Adverse Event Oncology Trial Response Team:
%   4-Site Rotation
%
% This file is a draft scaffold. A future Claude Code Opus 4.7 1M Max
% session reads the bracketed instruction blocks below, then reads the
% files those instructions name, then expands the scaffold into the
% finished section. Target length is roughly 6 to 8 pages of the 70
% plus page final paper.
% ======================================================================

\begin{sectionaccent}

[LIMITATIONS AND FUTURE WORK INSTRUCTIONS FOR FUTURE CLAUDE CODE OPUS
4.7 1M MAX SESSION. Write 5 subsections inside this single
\texttt{\textbackslash section} block. Every limitation must trace to a
file so the reader can verify the bound and so the future work
roadmap stays anchored to the real artifact.

Subsection 1, Simulation Versus Real Robot Limitations. Read:
\begin{itemize}
  \item \path{demo-projects/07-humanoid/paper/codegen/src/fda_rtct_submitter.py}
        (production uses ed25519 signing, simulation uses a placeholder
        signature).
  \item \path{demo-projects/07-humanoid/paper/codegen/src/robot_loop.cpp}
        (the in-robot 10 Hz loop runs in a simulator harness, not on
        the H2 EDU body).
  \item \path{demo-projects/07-humanoid/paper/codegen/src/sensor_to_xyz.py}
        (sensor frames are synthetic in the current run).
\end{itemize}
State that no real Unitree H2 EDU body executed any of the broadcast
commands in this paper. State that the FDA RTCT pilot intake API is
called only in simulation. State that the patient identifiers are
synthetic and have the form \texttt{PAT-NET-001-PNNN}.

Subsection 2, Per-Site Data Identity Limitation. Read the 4 site
files in
\path{demo-projects/07-humanoid/paper/execution/data/site_runtime_runs/}.
All 4 files share the same byte length (62,330 bytes each). State
that this parity is acceptable as a schema check and is not yet a
production behavior. State that the next iteration should feed each
site a site-specific AE arrival stream so the record content
diverges by site.

Subsection 3, Excluded Inputs. Read the Notes section of
\path{demo-projects/07-humanoid/paper/instructions/README.md}. The
original Prompt 07 extra hours dataset (hour 56 through hour 83 under
\path{physical-ai-oncology-trials/new-trial/national-24-7-trial/extra-hours/})
is excluded from the current input scope. Only hour 00 through hour
55 are read. State this explicitly and identify the missing hours as
a future inclusion roadmap item.

Subsection 4, Missing Repository Update Instruction Directory. Read
the Repository Structure block of
\path{demo-projects/07-humanoid/paper/instructions/README.md}. The
\path{10-repository-update-instructions/} subdirectory is named in
the README but is not yet present on disk. State this explicitly as a
documented limitation and as the next commit for the instruction
tree.

Subsection 5, Future Work Roadmap. Read:
\begin{itemize}
  \item \path{demo-projects/07-humanoid/paper/inputs/README.md}
        for the prior paper's future work framing.
  \item \path{demo-projects/07-humanoid/paper/inputs/Deep-Research-A/chunk2_tables_references.md}
        for the Medical Full Automation A roadmap items.
  \item \path{demo-projects/07-humanoid/paper/inputs/Deep-Research-B/chunk_04_comparative_overview_and_implications.md}
        for the Medical Rapid Prototyping B roadmap items.
  \item Author's December 2025 paper
        \path{kawchak_2025_18029100} for the prior limitations
        that this paper carries forward.
\end{itemize}
Roadmap items to enumerate explicitly:
\begin{itemize}
  \item Move from simulator harness to a real Unitree H2 EDU body
        with the kinematics in
        \path{codegen/src/kinematics.yaml} as the live target.
  \item Replace the placeholder signature in
        \path{codegen/src/fda_rtct_submitter.py} with the real
        ed25519 signing chain.
  \item Enroll a real PAT-NET-001 cohort under an IRB protocol and a
        signed sponsor agreement.
  \item Add the
        \path{instructions/10-repository-update-instructions/}
        directory.
  \item Extend the AE arrival stream so the 4 site decision logs
        diverge in content, not only in byte length.
  \item Include the hour 56 through hour 83 extra hours dataset from
        \texttt{physical-ai-oncology-trials} once the supporting
        analysis is ready.
  \item Upgrade the per-site compute path to NVIDIA Jetson AGX Thor
        2070 TOPS where the workload demands it.
\end{itemize}

Required real-life feasibility insight to include in this section:
``The largest gap between this paper and a real Phase 2 trial response
team is not technology. It is the IRB, sponsor, and FDA RTCT pilot
governance. The technology layer is feasible today. The governance
layer is the work that the paper invites the agencies, sponsors, and
sites to take on.''

Do not exceed 8 pages. Use single dashes only. Use \S\ for the
section sign. Avoid stranding 1 or 2 word lines on their own.]

The limitations and future work section will name the simulator
versus real robot gap, the per-site data identity limitation, the
excluded extra hours dataset, the missing repository update
instruction directory, and the future work roadmap including the IRB
and FDA RTCT pilot enrollment path.

\sectiontable{Limitations and Future Work planning matrix}{%
SimGap & L1 & 1 & robot, submit. & Sim vs real & 1 to 2 pages & No real H2 run & Ready \\
Parity & L2 & 2 & runtime runs & Identity check & 1 page & 62,330 bytes & Ready \\
Excluded & L3 & 3 & instr Notes & Hour 56 to 83 & 1 page & By README & Ready \\
Missing & L4 & 4 & instr structure & 10-repo gap & 1 page & Not present & Ready \\
Future & L5 & 5 & inputs, DR-A, DR-B & Roadmap list & 2 to 3 pages & 7 items & Ready \\
}

\end{sectionaccent}
